\newpage
\section{Random Matrix Theory}
In addition to non-asymptotic results, we will need asymptotic analysis, which is more delicate. The section is organized as follows: 
\begin{enumerate}
    \item Density of eigenvalues in classical ensembles of random matrices
    \item Semi-Circle Law and Marchenko–Pastur Law
    \item BBP Transition
    \item CLT for Eigenvalues
    \item Spectrum Separation
    \item Replica Method
\end{enumerate}

This section mainly follows STATC206B (UC Berkeley, taught by Vadim Gorin) and STAT260 (UC Berkeley, taught by Song Mei, 2021). I also refered to the book~\cite{bai2010spectral}.

\subsection{Density of Eigenvalues in Classical Ensembles of Random Matrices}

We begin our tour from an important type of matrix G$\beta$E ($\beta=1,2,4$). Let $X$ be $N\times N$ matrix with i.i.d. entries: 
\[
\begin{cases}
    a) \normalg \\
    b) \normalg + i\normalg \\
    c) \normalg + i\normalg + j\normalg + k\normalg
\end{cases}
\]
Set $M = \frac{1}{2}(X+X^*)$, and we have the main theorem.
\begin{thm}[Density of Eigenvalues]
    Let the eigenvalues of $M$ be $\lambda_1\ge \lambda_2\ge ...\ge \lambda_N$, and they have density
    \[
    \frac{1}{Z} \prod_{i<j}|\lambda_j-\lambda_i|^\beta \cdot \prod_{i=1}^N\exp(-\frac{\lambda_i^2}{2}) \dd \lambda_{1}...\dd \lambda_{N}.
    \]
    where $\beta = 1,2,4$ for a), b), c) and 
    \[
    Z = \frac{(2\pi)^{N/2}}{N!} \prod_{j=0}^N \frac{\Gamma(1+(j+1)\beta/2)}{\Gamma(1+\beta/2)}
    \]
    is the partition function (normalization factor).
\end{thm}
\begin{proof}
    We prove $\beta = 1$ here and leave $\beta = 2$ to the homework.

    \paragraph{Step 1}  Claim density of $M = \frac{1}{2}(X+X^*) \sim \exp(-\frac{1}{2}\tr(M^2))$.

    Indeed, 
    \[
    \tr(M^2) = \sum_{i,j} m_{ij}^2 = \sum_{i=1}^N \underbrace{x_{ii}^2}_{\mathcal{N}(0,1)} + \frac{1}{2}\sum_{i<j} \underbrace{(x_{ij}+x_{ji})^2}_{\mathcal{N}(0,2)}
    \]
    implies that the density of $M\propto \exp(-\frac{1}{2}\tr(M^2))$

    \paragraph{Step 2}  We can calculate that $\exp(-\frac{1}{2}\tr(M^2)) = \prod_{i=1}^N \exp(-\frac{\lambda_i^2}{2})$.

    We derive it immediately by diagonalizing the the matrix $M$.

    \paragraph{Step 3}  Each symmetric matrix is determined by its eigenvalues and eigenvectors, i.e. there exists an almost bijection $\pi$
    \[
    \pi: \underbrace{\mathcal{W}_N}_{\lambda_1<\lambda_2<...<\lambda_N} \times \underbrace{O(N)}_{\text{Orthogonal Bases}} \to \underbrace{\mathcal{H}_N}_{\text{Symmetric Matrix}}
    \]
    We say "almost" because indeed the map is not injective: we can multiply the eigenvalue by $\pm 1$. In other words, if the eigenvalues are unique, $\pi^{-1}(M)$ has exactly $2^N$ elements.

    Now we give the key proposition of the proof.
    \begin{prop}[Jacobian]
        Consider the map $\pi: (\Lambda, O) \mapsto B$, where $\pi((\Lambda, O)) = O\Lambda O^*$. Then the Jocobian of the map is $\prod_{i<j} |\lambda_j-\lambda_i|$.
    \end{prop}
    \begin{rem}
        It is an important technique to transform an intractable density calculation into a tractable one and a Jocobian.
    \end{rem}
    \begin{proof}
        It is sufficient to calculate by taking $O = \Id$, as when we transform $O$ to $A\cdot O$ (where $A$ is orthogonal), the \textbf{uniform measure} on $O(N)$ and the Lebesgue measure on $\mathcal{H}(N)$ remain unchanged.

        Then we take $O = \exp(B) = \Id + B + ...$ and we can deduce that $B + B^* = 0$ as $OO^* = \Id$. Then the map $\pi$ can be written as
        \begin{align*}
        \left((\lambda_1,...,\lambda_N), \exp(B)\right)
        &\mapsto \exp(B)
        \begin{pmatrix}
        \lambda_1 & 0 & 0 & \cdots & 0 \\
        0 & \lambda_2 & 0 & \cdots & 0 \\
        0 & 0 & \lambda_3 & \cdots & 0 \\
        \vdots & \vdots & \vdots & \ddots & \vdots \\
        0 & 0 & 0 & \cdots & \lambda_n
        \end{pmatrix}
        \exp(-B)\\
        &\mapsto (\Id+B)
        \begin{pmatrix}
        \lambda_1 & 0 & 0 & \cdots & 0 \\
        0 & \lambda_2 & 0 & \cdots & 0 \\
        0 & 0 & \lambda_3 & \cdots & 0 \\
        \vdots & \vdots & \vdots & \ddots & \vdots \\
        0 & 0 & 0 & \cdots & \lambda_n
        \end{pmatrix}
        (Id-B) + o(B)\\
        &\mapsto 
        \begin{pmatrix}
        \lambda_1 & 0 & 0 & \cdots & 0 \\
        0 & \lambda_2 & 0 & \cdots & 0 \\
        0 & 0 & \lambda_3 & \cdots & 0 \\
        \vdots & \vdots & \vdots & \ddots & \vdots \\
        0 & 0 & 0 & \cdots & \lambda_n
        \end{pmatrix}\\
        &+
        \begin{pmatrix}
        0 & b_{12}(\lambda_2 - \lambda_1) & b_{13}(\lambda_3 - \lambda_1) & \cdots & b_{1n}(\lambda_n - \lambda_1) \\
        b_{21}(\lambda_1 - \lambda_2) & 0 & b_{23}(\lambda_3 - \lambda_2) & \cdots & b_{2n}(\lambda_n - \lambda_2) \\
        b_{31}(\lambda_1 - \lambda_3) & b_{32}(\lambda_2 - \lambda_3) & 0 & \cdots & b_{3n}(\lambda_n - \lambda_3) \\
        \vdots & \vdots & \vdots & \ddots & \vdots \\
        b_{n1}(\lambda_1 - \lambda_n) & b_{n2}(\lambda_2 - \lambda_n) & b_{n3}(\lambda_3 - \lambda_n) & \cdots & 0
        \end{pmatrix}
        + o(B)
        \end{align*}
        As $B$ has only $\frac{n(n-1)}{2}$ free parameter, so the Jacobian of the map is $\prod_{i<j}|\lambda_j-\lambda_i|$.
    \end{proof}
    
    \paragraph{Step 4}  We now calculate $Z$:
    \begin{align*}
        Z = \int_{\mathbb{R}^N} \prod_{i<j}|\lambda_j-\lambda_i|^\beta \cdot \prod_{i=1}^N\exp(-\frac{\lambda_i^2}{2}) \dd \lambda_{1}...\dd \lambda_{N} = 
        \frac{(2\pi)^{N/2}}{N!} \prod_{j=0}^N \frac{\Gamma(1+(j+1)\beta/2)}{\Gamma(1+\beta/2)}
    \end{align*}
\end{proof}

\subsection{Semi-Circle Law and Marchenko–Pastur Law}
\subsubsection{Trace Calculation}
\begin{thm}[Tridiagonal Matrix Equivalent]
    Consider the real symmetric tridiagonal matrix
    \[
    T_{\beta} = 
    \begin{pmatrix}
    \cN(0,1) & \frac{1}{\sqrt{2}}\chi_{\beta(n-1)} & \cdots & 0 \\
    \frac{1}{\sqrt{2}}\chi_{\beta(n-1)} & \cN(0,1) & \cdots & 0 \\
    \vdots & \vdots & \ddots & \vdots \\
    0 & 0 & \cdots & \cN(0,1)
    \end{pmatrix}
    \]
    where the second diagonal entries are $\chi_{\beta(n-1)}, \chi_{\beta(n-2)},... \chi_{\beta}$.
    $\chi_k = \sqrt{\chi_k^2} = \sqrt{\sum_{i=1}^k \cN(0,1)^2}$, whose density is
    \[
    \frac{1}{2^{\frac{k}{2}-1}\Gamma(\frac{k}{2})}x^{k-1}e^{-\frac{x^2}{2}}.
    \]
    Then the eigenvalues of $T_\beta$ have the same distribution of $G\beta E$.
\end{thm}

\begin{proof}
    We present for $\beta = 1$. By linear algebra, we can choose an orthogonal matrix $U_1$, and transform the matrix $M$ to $U_1MU_1^\top$, whose the first row is 
    \[
    (x_{11}, \sqrt{\sum_{k=2}^n x_{ik}^2}, 0,...,0).
    \]
    Note that the eigenvalues do not change after the transformation.
    
    We can inductively do the same operation on the rest sub-matrix.
\end{proof}

\begin{cor}
    For $\beta = 1, 2, 4$, the law of eigenvalues $\sim \prod_{i<j} (x_i-x_j)^\beta \prod_{i=1}^N\exp(-\frac{x_i^2}{2})$.
\end{cor}

\begin{thm}
    In fact, the corollary is true for any $\beta>0$.
\end{thm}

Now we begin to prove the semi-circle law. First we give the key result of trace calculation.
\begin{thm}[Semi-Circle Law]
    For each $k>0, \beta>0$, we have
    \[
    \frac{1}{N^{\frac{k}{2}+1}} \tr(T_\beta^k)=
    \frac{1}{N^{\frac{k}{2}+1}}\sum_{i=1}^N X_i^k
    \xrightarrow{N\to \infty, \text{in probability}} 
    \begin{cases}
        0&, \quad \text{k is odd}; \\
        (\frac{\beta}{2})^k \Cat_{\frac{k}{2}}&, 
        \quad \text{k is even}.
    \end{cases}
    \]
    where $\Cat_k$ is the k-the catalan number defined as
    \[
    \Cat_k = \frac{1}{k+1} \binom{2k}{k}.
    \]
\end{thm}

\begin{proof}
    By SLLN, we have
    \begin{align*}
        &\lim_{N\to \infty} \frac{\chi_{\beta(N-1)}}{\sqrt{2N}} = \sqrt{\frac{\beta}{2}},\\
        &\lim_{N\to \infty} \frac{\chi_{\beta(N-\alpha N)}}{\sqrt{2N}} = \sqrt{\frac{\beta(1-\alpha)}{2}}.
    \end{align*}
    So by the definition of the trace, we have
    \begin{align*}
    \tr \ \left(\frac{T_\beta}{\sqrt{N}}\right)^k 
    &= \sum_{m=1}^N \sum_{\substack{\text{k-step paths} \\ \text{from } m \text{ to } m}} 
    \left(\ \textcolor{gray}{ \text{some} 
    \ \frac{\chi_{\beta(N-i)}}{\sqrt{2N}} } \right)
    \left(\ \textcolor{gray}{ \text{some} 
    \ \frac{\mathcal{N}(0,1)}{\sqrt{N}} }\right)\\
    &=\sum_{m=1}^N \left( \# {\text{k-step paths from } m \text{ to } m} \right) \sqrt{\frac{\beta}{2}\left( \frac{N-m}{N} \right)^k} + o(N).
    \end{align*}
    The second equation follows from that if there exists "some $\frac{\mathcal{N}(0,1)}{\sqrt{N}}$", the term vanishes. 
    When $2 \nmid k$, $\# {\text{k-step paths from } m \text{ to } m}$ is zero.
    When $2\mid k$, $\# {\text{k-step paths from } m \text{ to } m} = \binom{k}{\frac{k}{2}}$. Approximate the summation by integration and we have
    \[
    \frac{1}{N} \tr \ \left(\frac{T_\beta}{\sqrt{N}}\right)^k \to \binom{k}{\frac{k}{2}} \left(\frac{\beta}{2}\right)^{\frac{k}{2}} \int_0^1 x^{\frac{k}{2}}\dd x = \Cat_{\frac{k}{2}} \left(\frac{\beta}{2}\right)^{\frac{k}{2}}
    \]
\end{proof}

Now we introduce the semi-circle distribution.
\begin{defn}[Wiegner Semi-Circle Distribution]
    The density of the distribution is
    \[
    \mu(x) = \frac{1}{2\pi} \sqrt{4-x^2}.
    \]
\end{defn}

We use moment method to recover the proof and give the definition.
\begin{cor}[Moment Analysis]
    Let $m_k$ are moments of $T_\beta$ and we have
    \[
    m_u = \begin{cases}
        0&,\quad u \text{ is odd}; \\
        \frac{1}{\frac{u}{2}+1}\binom{u}{\frac{u}{2}} &,\quad u \text{ is even}.
    \end{cases}
    \]
\end{cor}

\begin{thm}[Moment of Semi-Circle Distribution]
    \label{thm:scl}
    $m_k$ are moments of semi-circle law and
    \[
    m_k = \int_{-2}^2 \mu(x) x^k \dd x.
    \]
\end{thm}

\begin{proof}{(Method I)}
    We just calculate
    \[
    m_k = \int \frac{1}{\sqrt{2\pi}} \sqrt{4-x^2}x^k \dd x.
    \]
    Let $x = 2\cos\theta$, and by inductive calculation, we derive the results.
\end{proof}

However, we are not satisfied for the calculation is not so intuitive. We then introduce another proof.

\begin{proof}{(Method II)}
Introduce generating function: 
\[
G(z) = \sum_{k=0}^\infty m_k z^{-k-1},
\]
$m_0 =1$. By classical results, we have
\[
\sum_{n=0}^\infty \Cat_n x^n = \frac{1-\sqrt{1-4x}}{2x} =: C(x).
\]
Now we can derive the expression of $G(z)$ using $C(x)$.
\[
    G(z) = \sum_{k=0}^\infty m_k z^{-k-1} = \sum_{l=0}^\infty \Cat_l z^{-2l-1} = \frac{z-\sqrt{z^2-4}}{2}.
\]
Next we introduce two propositions.
\begin{prop}[Stieltjes Transform]
    \label{prop:2.8}
    \[
    G(z) = \int \frac{\mu(x)}{z-x} \dd x.
    \]
\end{prop}
\begin{proof}
    By Taylor expansion, we have
    \begin{align*}
    \int \frac{\mu(x)}{z-x} \dd x 
    &= \frac{1}{z} \int \frac{\mu(x)}{1-\frac{x}{z}}\dd x\\
    &= \frac{1}{z}\int \mu(x)\sum_{k=0}^\infty\left( \frac{x}{z}\right)^k\\
    &=\sum_{k=0}^\infty m_k z^{-k-1} \\
    &= G(z)
    \end{align*}
\end{proof}

\begin{prop}
    \label{prop:2.9}
    \[
    \mu(x_0) = -\frac{1}{\pi} \lim_{y_0\to 0^+} \Im (G(x_0+iy_0)).
    \]
\end{prop}
\begin{rem}
    We can just add a perturbation onto the imagine axis and obtain the information of the point.
\end{rem}

\begin{proof}
    By proposition \ref{prop:2.8}, we have
    \begin{align}
        -\frac{1}{\pi}  \Im (G(x_0+iy_0))
        &= -\frac{1}{\pi} \int  \Im \frac{1}{x_0+iy_0-x} \mu(x)\dd x\\
        &= -\frac{1}{\pi} \int \Im \frac{x_0-x-iy_0}{(x-x_0)^2+y_0^2}\mu(x) \dd x\\
        &= \int \frac{1}{\pi} \frac{y_0}{(x-x_0)^2+y_0^2}\mu(x) \dd x.
    \end{align}
    Notice that $\frac{1}{\pi} \frac{y_0}{(x-x_0)^2+y_0^2}$ is a "good" kernel, and thus as $y_0\to 0$
    \[
    \int \frac{1}{\pi} \frac{y_0}{(x-x_0)^2+y_0^2}\mu(x) \dd x \to \mu(x_0).
    \]
\end{proof}
Return to the Theorem \ref{thm:scl}, by proposition \ref{prop:2.9}
\[
\mu(x_0) = -\frac{1}{\pi} \lim_{y\to 0^+} \Im \left( \frac{1}{2}\left[(x+iy) - \sqrt{(x+iy)^2 - 4)}\right]
\right) = \begin{cases}
    0  &, \quad |x|>2;\\
    \frac{\sqrt{4-x^2}}{2\pi} &,\quad |x|\le 2.
\end{cases}
\]
\end{proof}

Now we are prepared to formally state and prove the semi-circle law.
\begin{thm}[Semi-Circle Law]
    Let $\lambda_1 < ... <\lambda_N$ be eigenvalues of G$\beta$E. Set $x_i = \lambda_i\sqrt{\frac{2}{\beta N}}$ and let $\mu_N$ be their empirical measure: $\mu_N = \frac{1}{\sqrt{N}}\sum_{i=1}^N \delta_{x_i}$. Then
    \[
    \lim_{N\to \infty} \mu_N = \frac{1}{2\pi} \sqrt{4-x^2} \mathbbm{1}_{|x|\le 2} = \mu_{\text{circle}}(x)
    \]
    weakly in probability, which means $\forall$ bounded continuous function $f(x)$
    \[
    \lim_{N\to \infty} \int_{\RR} f(x)\mu_N(x) = \int_{\RR} f(x)\mu(x) \dd x.
    \]
\end{thm}

\begin{proof}
    We prove the result in four steps.
    \paragraph{Step 1}  The results hold for $f(x) = x^k$.
    
    \paragraph{Step 2}  The results hold for any polynomials. 
    
    \paragraph{Step 3}  Take $L>2$, the results hold for $f(x) = \mathbbm{1}_{|x|>L} x^k$.
    We have
    \begin{align*}
        \left|\int f(x)\mu_N(\dd x)\right| 
        &= \left|\int x^k\mu_N(\dd x)\right|\\
        &\le L^{-2m} \int_{|x|\ge L} x^{2k+2m}\mu_N(\dd x)\\
        &\to L^{-2m} \int x^{2k+2m} \mu(\dd x)\\
        &\le L^{-2m} 2^{2k+2m} = 2^{2k} \left(\frac{2}{L}\right)^m \to 0. 
    \end{align*}
    
    \paragraph{Step 4}  By step 3, we can restrict the support set of $f$ on a compact set, i.e. $[-4, 4]$. Apply the Weierstrass theorem and we get the proof.
\end{proof}

At last, we give three generalizations of the semi-circle law. Next theorem tells us the gaussian assumption of the semi-circle law is not neccessary.
\begin{thm}[Generalization]
    Let $z_{ij}, i<j$ be i.i.d. random variables with finite moments. $\EE z_{ij} = 0$ and $\EE z_{ij}^2 = \frac{1}{2}$. Let $Y_i$ be i.i.d. with finite moments. Then the semi-circle law still holds for the matrix with entries $Y_i$ and $z_{ij}$.
\end{thm}

\subsubsection{Marchenko–Pastur Law}
Now we consider the case where the matrix is not square.
Let $X$ be an $N\times S$ random matrix whose elements are Gaussian with parameter $\beta > 0$ (different form according to $\beta$).
Assume that $N, S\to\infty$ in such a way that $N/S \to \gamma^2 \in (0,1)$.
Let $\lambda_1, \lambda_2,\ldots, \lambda_N$ denote the eigenvalues of $\frac{1}{S}XX^*$, and define the empirical spectral measure
\[
\rho^N = \frac{1}{N}\sum_{i=1}^N \delta_{\lambda_i}.
\]
 
\begin{thm}[Marchenko--Pastur Law]
As $N\to\infty$ (with $N/S\to \gamma^2\in(0,1)$), the empirical spectral measure $\rho^N$ converges weakly, in probability, to the Marchenko--Pastur distribution $\mu_{\mathrm{MP}}$ with density
\[
\frac{\dd\mu_{\mathrm{MP}}}{\dd x}(x)
= \frac{1}{2\pi\beta\gamma^2}\,
\frac{\sqrt{(\lambda_+ - x)(x-\lambda_-)}}{x}\,
\mathbf{1}_{[\lambda_-,\,\lambda_+]}(x),
\]
where
\[
\lambda_+ = \beta(1+\gamma)^2, \qquad
\lambda_- = \beta(1-\gamma)^2.
\]
\end{thm}

\begin{proof} We calculate the moments for $\rho^N$ and $\mu_{\mathrm{MP}}$ separately
    \paragraph{Step 1} We claim that
    \[
    \lim_{N,S\to \infty} \frac{1}{NS^k}\tr[(XX^*)^k] = \beta^k \sum_{r=0}^{k-1} \frac{\gamma^{2r}}{r+1} \binom{k}{r} \binom{k-1}{r}.
    \]

    Notice that as $S\to \infty$
    \[
    \frac{\chi_{aS}}{\sqrt{S}} \to \sqrt{a},
    \]
    we have
    \begin{align*}
        \operatorname{LHS} &= \lim_{S,N\to\infty}\frac{1}{NS^k} \sum_{m=1}^N\sum_{\substack{\text{k-step paths} \\ \text{from } m \text{ to } m}} (\text{diagonal entries})(\text{sub-diagonal entries})\\
        &= \lim_{S,N\to \infty} \frac{1}{NS^k} \sum_{m=1}^N \sum_{i=0}^{[k/2]} \left(
        \chi_{\beta(N-m)}\chi_{\beta(S-m+1)}
        \right)^{2i}\chi_{S+N-2m-2}^{2(k-2i)} \binom{k}{2i}\binom{2i}{i}
    \end{align*}
    By replacing $\chi$ by its limit, we have
    \[
    \frac{\chi_{\beta(N-m)}\chi_{\beta(S-m+1)}}{S} \to \beta \sqrt{(\gamma^2-\frac{m}{S})(1-\frac{m}{S})}
    \]
    \[
    \frac{\chi^2_{\beta(S+N-2m+2)}}{S} \to \beta(1+\gamma^2 - 2\frac{m}{S})
    \]
    Then we approximate the sum by integral
    \begin{align*}
    \operatorname{LHS} 
    &= \frac{\beta^k}{\gamma^2} \sum_{i=0}^{[k/2]} \binom{k}{2i}\binom{2i}{i} \int_0^{\gamma^2} (\gamma^2-t)^i(1-t)^i (1+\gamma^2-2t)^{k-2i}\dd t\\
    &=\frac{\beta^k}{\gamma^2} \sum_{i=0}^{[k/2]} \binom{k}{2i}\binom{2i}{i} \int_0^{\gamma^2} u^i(1-\gamma^2+u)^i (1-\gamma^2+2u)^{k-2i}\dd u
    \end{align*}
    We claim that
    \[
    \frac{1}{\gamma^2}\sum_{i=0}^{[k/2]} \binom{k}{2i}\binom{2i}{i} \int_0^{\gamma^2} u^i(1-\gamma^2+u)^i (1-\gamma^2+2u)^{k-2i}\dd u = \sum_{r=0}^{k-1} \frac{\gamma^{2r}}{r+1} \binom{k}{r} \binom{k-1}{r}.
    \]
    and we obtain the proof.

    The proof of the claim: 
    Denote the left hand side as A. First, we observe that
    \begin{equation}
        \label{eq:key}
        \operatorname{A} = \text{The sum of coefficients of all the }x^iy^i \text{ in} \int_{0}^{\gamma^2} [xu+y(1-\gamma^2 +u)+(1-\gamma^2+2u)]^k\dd u
    \end{equation}
    Denote B as the integral of (\ref{eq:key})
    \begin{align*}
        B &= \frac{1}{\gamma^2}\frac{[(x+1)\gamma^2+(y+1)]^{k+1} - [(y+1)(1-\gamma^2)]^{k+1}}{(x+y+2)(k+1)}\\
        &=\frac{1}{k+1}\sum_{l=0}^k [(y+1)(1-\gamma^2)]^l [(x+1)\gamma^2+y+1]^{k-l}.
    \end{align*}
    To calculate the coefficients, we can let $x = \frac{1}{y}$ and calculate the constant coefficient.
    \begin{align*}
        \operatorname{B} &= \frac{1}{k+1} \sum_{l=0}^k [(y+1) (1-\gamma^2)]^l
        \left(\left(\frac{1}{y}+1\right)\gamma^2 + y+1\right)^{k-l}\\
        &=\frac{(y+1)^k}{k+1}
        \sum_{l=0}^k(1-\gamma^2)^l(\frac{1}{y}\gamma^2+1)^{k-l}\\
        &=\frac{(y+1)^k}{k+1} \frac{(\frac{1}{y}\gamma^2+1)^{k+1} - (1-\gamma^2)^{k+1}}
        {\frac{1}{y}\gamma^2 + \gamma^2}\\
        &=\frac{y(y+1)^{k-1}}{(k+1)\gamma^2} \left(\left(\frac{1}{y}\gamma^2+1\right)^{k+1} - (1-\gamma^2)^{k+1}\right)
    \end{align*}
    Then the constant coefficient is
    \[
    \frac{1}{(k+1)\gamma^2} \sum_{r=0}^{k-1} \binom{k-1}{r} \binom{k+1}{r+1} \gamma^{2r+2} = \sum_{r=0}^{k-1} \frac{\gamma^{2r}}{r+1} \binom{k}{r} \binom{k-1}{r}.
    \]
    
    \paragraph{Step 2} We claim that $\lambda_{+} = \beta(1+\gamma)^2$ and $\lambda_{-} = \beta(1-\gamma)^2$.
    
    Define the density of the limit distribution
    \[
    \mu(x) = \frac{1}{2\pi\gamma^2}\frac{\sqrt{(\lambda_{+}-x)(x-\lambda_{-})}}{x} \mathbbm{1}_{[\lambda_{-}, \lambda_{+}]}
    \]
    and $\mu_N = \rho^N$.
    
    The calculation below is inspired by \href{http://ndl.ethernet.edu.et/bitstream/123456789/60805/1/140.pdf}{Zhidong Bai's book}.
    First, we calculate the moment of $\mu(x)$, for $m\ge 1$,
    \begin{align*}
        \EE_{\mu} x^m &= \int x^m \mu(x)\dd x\\
        &= \frac{1}{2\pi\gamma^2}\int x^{m-1} \sqrt{(\lambda_{+}-x)(x-\lambda_{-})} \dd x\\
        &= \frac{1}{2\pi\gamma^2}(2\gamma)^2\int_{-1}^1 (1+\gamma^2+2\gamma u)^{m-1}\sqrt{1-u^2}\dd u
        \\
        &(\text{by setting }x =\beta(1+\gamma^2+2\gamma u))\\
        &= \beta^m\frac{1}{\pi}\sum_{l=0}^{[(m-1)/2]} \binom{m-1}{2l}(1+\gamma^2)^{m-1-2l}
        (4\gamma^2)^l
        \int_{-1}^1 u^{2l}\sqrt{1-u^2}\dd u\\
        &= \beta^m
        \sum_{l=0}^{[(m-1)/2]} \binom{m-1}{2l} 
        \frac{1}{l+1}\binom{2l}{l} (1+\gamma^2)^{m-1-2l} \gamma^{2l}\\
        &= \beta^m\sum_{l=0}^{[(m-1)/2]}\sum_{s=0}^{m-1-l} \frac{(m-1)!}{l!(l+1)!s!(m-1-2l-s)!}\gamma^{2(l+s)}\\
        &=\beta^m \frac{1}{m}\sum_{r=0}^{m-1} \gamma^{2r} \binom{m}{r} \sum_{l=0}^{r\wedge(m-1-r)}\binom{r}{l} \binom{m-r}{m-r-l-1}\\
        &=\beta^m \frac{1}{m}\sum_{r=0}^{m-1}\binom{m}{r}\binom{m}{r+1}\gamma^{2r}\\
        &=\beta^m \sum_{r=0}^{m-1}\frac{1}{r+1}\binom{m}{r}\binom{m-1}{r}\gamma^{2r}.
    \end{align*}

    Next, we notice that
    \[
        \EE_{\mu_N} x^m =\frac{1}{N} \sum_{i=1}^N \lambda_i^N = \frac{1}{N}\tr\left[\left(\frac{XX^*}{S}\right)^m\right]
    \]
    and
    \[
    \lim_{N\to \infty} \frac{1}{N}\tr\left[\left(\frac{XX^*}{S}\right)^m\right] = \beta^m \sum_{r=0}^{m-1}\frac{1}{r+1}\binom{m}{r}\binom{m-1}{r}\gamma^{2r}.
    \]
    So we have
    \[
    \lim_{N\to \infty}\EE_{\mu_N} x^m = \EE_{\mu} x^m.
    \]
    Similar to the proof of semi-circle law, we have $\rho^N$ converges (weakly, in probability) to the Marchenko-Pastur law.
\end{proof}


\subsection{BBP Transition}
In the spiked model $\Sigma = I + \theta\, vv^*$, Baik, Ben Arous, and P\'{e}ch\'{e} (2005) discovered a phase transition for $\lambda_{\max}$ of $\Sigma$. In general form, it is stated as

\begin{thm}[Spiked Random Matrices]
    \label{thm:9.1}
    Let $\lambda_1<\lambda_2<\cdots<\lambda_N$ are $N$ eigenvalues of 
    \[
    C = \sqrt{N}a vv^* +\sqrt{\frac{2}{\beta}}G\beta E,
    \]
    where $\|v\|=1$
    Then there exists $a_{crit}$ such that
    \begin{enumerate}
        \item If $a>a_{crit} = 1$, then 
        \[
        \lim_{N\to\infty} \frac{\lambda_N}{\sqrt{N}} = a+\frac{1}{a} >2.
        \]
        \item If $a<a_{crit}$, then 
        \[
        \lim_{N\to\infty} \frac{\lambda_N}{\sqrt{N}} = 2.
        \]
    \end{enumerate}
\end{thm}
By the theorem, we can easily get
\begin{cor}
    \[
    \hat{a} = \frac{1}{2}\left(
    \frac{\lambda_N}{\sqrt{N}}+\sqrt{\frac{\lambda_N^2}{N}-4}
    \right) \overset{N\to\infty}{\longrightarrow} a.
    \]
\end{cor}
\begin{thm}[Recover $v$]
    \label{thm:9.2}
    In the setting of Theorem \ref{thm:9.1}. Let $\hat{v}$ denote the eigenvector corresponding to $\lambda_N$ and let $\phi$ be the angle between $v$ and $\hat{v}$,
    \[
    \lim_{N\to\infty} \sin\phi = \frac{1}{a} \wedge 1.
    \]
\end{thm}
\begin{proof}
    We prove the results in three steps.
    \paragraph{Step 1} \ $a, \lambda_N$ and $\phi$ are unchanged under orthogonal/unitary transformations of matrix $C$. We do two transformations:
    \begin{itemize}
        \item Rotate $v$ ti be the first basis vector $(1, 0^{N-1})$.
        \item Rotate in the orthogonal complement of $(1, 0^{N-1})$, so that the $(N-1)\times (N-1)$ bottom random corner of G$\beta$E becomes diagonal, while the law of the first row and column of G$\beta$E preserved.
    \end{itemize}
    \paragraph{Step 2} \ Now we brought $C$ into
    \[
    \hat{C} =
    \begin{pmatrix}
    a \delta_N + \mathcal{N}(0, \frac{2}{\beta}) & \xi_2 & \xi_3 & \dots & \xi_N \\
    \bar{\xi}_2 & \mu_2 &  &  &  \\
    \bar{\xi}_3 &  & \mu_3 &  &  \\
    \vdots &  &  & \ddots &  \\
    \bar{\xi}_N &  &  &  & \mu_N
    \end{pmatrix}
    \]
    and 
    $\mu_2,\cdots,\mu_N$ be the eigenvalues of $\sqrt{\frac{2}{\beta}} G\beta E$ and $\xi_2,\cdots,\xi_N \overset{i.i.d}{\sim} \sqrt{\frac{2}{\beta}}$ corresponding Normal distribution.
    We find an eigenvector $(x_1,\cdots,x_N)$ of $\tilde{C}$ with eigenvalues $\lambda$
    \[
    \begin{cases}
        x_1\left(a\sqrt{N}+\cN\left(0,\frac{2}{\beta}\right)\right)+\sum_{i=1}^N\xi_ix_i=\lambda x_1\\
        x_1\bar{\xi}_2+\mu_2x_2=\lambda x_2\Rightarrow x_2=\frac{x_1\bar{\xi}_2}{\lambda-\mu_2}\\
        \qquad\qquad\qquad\vdots\\
        x_1\bar{\xi}_N+\mu_Nx_N=\lambda x_N\Rightarrow x_N=\frac{x_1\bar{\xi}_N}{\lambda-\mu_N}
    \end{cases}
    \Longrightarrow
    \begin{cases}
        x_2 = \frac{\bar{\xi}_2}{\lambda-\mu_2}\\
        \qquad\vdots\\
        x_N = \frac{\bar{\xi}_N}{\lambda-\mu_N}
    \end{cases}
    \]
    Plug the $2 \sim N$ equations into the first line and we get
    \begin{equation}
        x_1 \left[a\sqrt{N}+\cN\left(0,\frac{2}{\beta}\right) - \lambda + \sum_{i=1}^N\frac{\xi_i\bar{\xi}_i}{\lambda-\mu_i}
        \right] = 0. \tag{$\ast$}
    \end{equation}
    By interpolating of eigenvalues, $\title{C}$ has 1 eigenvalue larger than $\mu_N$. Only this eigenvalue has chance to become larger than $2\sqrt{N}$ as $\frac{\mu_N}{\sqrt{N}} \to 2$. Denote $y = \frac{\lambda_N}{\sqrt{N}}$ and investigate ($\ast$) for $y>2$.
    \[
    a + \frac{\cN\left(0,\frac{2}{\beta}\right)}{\sqrt{N}} - y + \frac{1}{N} \sum_{i=2}^N \frac{\xi_i\bar{\xi}_i}{y - \frac{\mu_i}{\sqrt{N}}} = 0. \tag{$\ast\ast$}
    \]
    By LLN, we have
    \[
    \frac{1}{N} \sum_{i=2}^N \frac{\xi_i\bar{\xi}_i}{y - \frac{\mu_i}{\sqrt{N}}} \approx \frac{1}{N} \sum_{i=2}^N \frac{1}{y - \frac{\mu_i}{\sqrt{N}}} \to G(y) = \frac{1}{2}(y- \sqrt{y^2-4}).
    \]
    Then we have
    \[
    a - y + \frac{1}{2}(y- \sqrt{y^2-4}) = 0 \Longrightarrow y = a+ \frac{1}{a},
    \]
    and this proves Theorem \ref{thm:9.1}.

    \paragraph{Step 3} \ For Theorem \ref{thm:9.2}, we notice that the eigenvalues are
    \[
    (1,0,\cdots,0) \text{ and } \left(1, \frac{\bar{\xi}_2}{\lambda-\mu_2},\cdots,\frac{\bar{\xi}_N}{\lambda-\mu_N}
    \right),
    \]
    and we have
    \[
    \cos^2 \phi = \frac{1}{1+\sum_{i=2}^N\frac{\xi_i\bar{\xi}_i}{(\lambda-\mu_i)^2}} \to \frac{1}{1+\frac{1}{2}\frac{y}{\sqrt{y^2-4}}-\frac{1}{2}} = \frac{2\sqrt{y^2-4}}{y+\sqrt{y^2-4}}.
    \]
    As $y = a+\frac{1}{a}$, we have
    \[
    \sin^2 \phi \to \frac{1}{a^2}.
    \]
\end{proof}

\subsection{CLT for Eigenvalues}
Similar to the CLT for random variables, we have the following theorem.
\begin{thm}
\label{RMT_CLT}
Let $f$ be analytic in a small neighboring of $[-2,2]$. Let $\lambda_1\le\cdots\le\lambda_N$ be e.v. of $\sqrt{\frac{2}{\beta}}G\beta E,\beta>0$. Then
\[
\sum_{i=1}^N f\left(\frac{\lambda_i}{\sqrt{N}}\right)-N\int f(x)\frac{1}{2\pi}\sqrt{4-x^2}
\]
converges to dist to a Gaussian r.v. $\xi_f$ jointly other several $f=f_1,\dots,f_K$ with
\[
\mathbb{E}[\xi_f]=\left(\frac{2}{\beta}-1\right)m(f),\quad \operatorname{Cov}(\xi_f,\xi_g)=\frac{2}{\beta}C(f,g)
\]
\[
m(f)=\frac{1}{4}(f(2)+f(-2))-\frac{1}{2\pi}\int_{-2}^2
\frac{f(x)}{\sqrt{4-x^2}}\dd x
\]
\[
C(f,g) = \frac{1}{4\pi} \int_{-2}^2 \int_{-2}^2 \frac{(f(x)-f(y))(g(x)-g(y))}{(x-y)^2}\frac{4-xy}{\sqrt{4-x^2}\sqrt{4-y^2}}\dd x\dd y.
\]
or
\[
C\left(\frac{1}{z-x},\frac{1}{w-x}\right)=-\frac{1}{2(z-w)^2}\left(1-\frac{zw-4}{\sqrt{z^2-4}\sqrt{w^2-4}}\right)
\quad \text{for $z,w$ out of }[-2,2]
\]
where $m(f)$ and $C(f,g)$ do not depend on $\beta$.
\end{thm}

\begin{proof}
    The main idea is fancy moments method.

\begin{lem}[Wick's formula]
\label{lem:12.1}
\[
\mathbb{E}[\eta_1,\dots,\eta_m]=\sum_{\operatorname{perfect matchings}}\prod_{(i,j)\in\operatorname{mathching}}\sigma_{ij}
\]
\end{lem}
\begin{proof}
    By Laplacian's transform
    \[
    \EE [e^{t_1\eta_1+\cdots +t_m\eta_m}] = \exp\left(\frac{1}{2} \sum_{i,j=1}^m t_it_j\sigma_{ij} \right).
    \]
    Hence, we take derivative and get the expectation
    \[
    \EE [\eta_1\cdots\eta_m] = \frac{\partial^m}{\partial t_1\cdots\partial t_m} \left[\exp\left(\frac{1}{2} \sum_{i,j=1}^m t_it_j\sigma_{ij} \right)\right] \bigg|_{t_1 = \cdots = t_m = 0}.
    \]
    Equivalently, this is the coefficient of $t_1\cdots t_m$ in Taylor's expansion.
    \[
    = \sum_{ij} 2^{-\frac{m}{2}} \prod_{ij} \sigma_{ij}.
    \]
    Notice that each perfect match appears $2^{\frac{m}{2}}$ times, we prove the Lemma \ref{lem:12.1}.
\end{proof}

\paragraph{Takeaway} Moments can be reconstructed by applying a differential operator to Laplace transform.

\begin{lem}
\label{lem:12.2}
Introduce an operator
\[
\mathcal{D}_a=\prod_{i,j}(z_i-z_j)^{-1}\left(\sum_{i=1}^NT_{a,i}\right)\prod_{i,j}(z_i-z_j)
\]
\[
T_{a,i}f(z_1,\dots,z_N)=f(z_1,\dots,z_{i-1},z_i+a,z_{i+1},\dots,z_N)
\]
Then for $\lambda_1,\cdots,\lambda_N$ eigenvalues of GUE, we have 
\[
\underbrace{\EE \prod_{k=1}^M \left[\sum_{i=1}^N e^{a_k\lambda_i}\right]}_{\text{Moments of linear statistic for $f(\lambda)=e^{a\lambda}$}} = \underbrace{\mathcal{D}_{a_m} \cdots \mathcal{D}_{a_1} }_{\text{They all commute}} \exp\left(\sum_{i=1}^N \frac{z_i^2}{2}\right) \bigg|_{z_1 = \cdots = z_N = 0}.
\]
\end{lem}
\begin{proof}
From Lecture 6,
\[
\EE\exp(\operatorname{Tr(GUE\cdot Z)})=\EE B_{\lambda_1,\dots,\lambda_N}(z_1,\dots,z_N)=\exp\left(\sum_{i=1}^N\frac{\lambda_i^2}{2}\right)
\]
where $Z=\operatorname{diag}(z_1,\dots,z_N)$ and $\lambda_1,\dots,\lambda_N$ denote the e.v. of GUE.

Then act with $\mathcal{D}$
\[
B_{\lambda_1,\dots,\lambda_N}(z_1,\dots,z_N)=\prod_{K=1}^N(K!)\cdot\frac{\operatorname{det}[\exp(\lambda_iz_j)]}{\prod_{i,j}(\lambda_i-\lambda_j)(z_i-z_j)}.
\]
\[
\mathcal{D}=\prod_{i,j}(z_i-z_j)^{-1}\left(\sum_{i=1}^NT_{a,i}\right)\prod_{i,j}(z_i-z_j).
\]
So $B$ is eigenfunction of $\mathcal{D}$ with eigenvalue $\sum_{i=1}^N \exp(a\lambda_i)$. Hence,
\[
\EE\left[
\prod_{k=1}^m \left(\sum_{i=1}^N e^{a_k\lambda_i}\right)
\right] B_{\lambda_1 \cdots \lambda_N} (z_1,\cdots,z_N) = \mathcal{D}_{a_m} \cdots \mathcal{D}_{a_1} \exp\left(\sum_{i=1}^N\frac{\lambda_i^2}{2}\right).
\]
Plug $z_1,\dots,z_N=0$ and notice $B(0,\cdots,0)=1$.
\end{proof}
\begin{lem}
\label{lem:12.3}
\[
\mathcal{D} f(z_1)\cdots f(z_N)= f(z_1)\cdots f(z_N)\frac{a^{-1}}{2\pi i}\oint_{\{z_1,\dots,z_N\}}\left[\prod_{j=1}^N\frac{v+a-z_j}{v-z_j}\right]\frac{f(v+a)}{f(v)}
\]

The contour encloses $z_1,\dots,z_N$ with no singularities of $f$.
\end{lem}
\begin{proof}
\begin{align*}
    \mathcal{D} f(z_1) \cdots f(z_N) &= \sum_{i=1}^N \left[\prod_{j\neq i} \frac{z_i+a-z_j}{z_i-z_j}\right]\frac{f(z_i+a)}{f(z_i)}\cdot f(z_1)\cdots f(z_N)\\
    &=\frac{a^{-1}}{2\pi i}\oint_{\{z_1,\dots,z_N\}}\prod_{i=1}^N\frac{v+a-z_j}{v-z_j}\cdot\frac{f(v+a)}{f(v)}\dd v f(z_1)\cdots f(z_N)
\end{align*}
\end{proof}

\begin{cor}
\label{RMT_CLT_cor}
    \begin{align*}
    \EE\left[\prod_{K=1}^m\sum_{i=1}^N e^{a_K\lambda_i}\right] = &\frac{(a_1\cdots a_m)^{-1}}{(2\pi i)^m} \oint \sum_{k=1}^m \left[ \frac{(v_k+a_k)^N}{v_k^N} \exp\left(\frac{a_k^2}{2} + a_kv_k\right)\right] \\
    &\prod_{k<l} \frac{v_k-v_l+a_k-a_l}{v_k-v_l-a_l} \cdot \frac{v_k-v_l}{v_k-v_l+a_k} \dd v_1\cdots \dd v_k.
    \end{align*}
\end{cor}
\begin{proof}
We compute $\mathcal{D}_{a_m} \cdots \mathcal{D}_{a_1} \exp(\frac{z_i^2}{2})\cdots\exp(\frac{z_N^2}{2})$. By sequentially applying Lemma \ref{lem:12.3} and then setting $z_1=\cdots=z_N=0$ at the end
\end{proof}

Now we prove Theorem~\ref{RMT_CLT} for $\beta=2,f(\frac{\lambda}{\sqrt{N}})=\exp(a\cdot\frac{\lambda}{\sqrt{N}})$

\paragraph{Step 1: Expectation}  By Corollary~\ref{RMT_CLT_cor} with $m=1$
\[
\EE\left[\sum\exp\left(a\cdot\frac{\lambda_i}{\sqrt{N}}\right)\right]=\frac{(\frac{a}{\sqrt{n}})^{-1}}{2\pi i}\oint_{\{0\}}\left(\frac{v+\frac{a}{\sqrt{N}}}{v}\right)^N\exp\left(\frac{a^2}{2N}+\frac{a}{\sqrt{N}}v\right)\dd v
\]
No steepest descent needed. Set $v=u\sqrt{N}$ to get
\begin{align*}
    &\frac{N}{a} \frac{1}{2\pi i} \oint \exp\left( N\log\left(1+\frac{a}{Nu}\right) + au + \frac{2a^2}{N}\right)\\
    =& \frac{N}{a} \frac{1}{2\pi i} \oint \exp\left(a\left(u+\frac{1}{u}\right) + \frac{a^2}{2N} \left( 1-\frac{1}{u^2} \right)+O(N^{-2})\right)\dd u\\
    =& \frac{N}{a} \frac{1}{2\pi i} \oint \exp\left(a\left(u+\frac{1}{u}\right)\right) \dd u + \frac{a}{4\pi i} \oint \exp\left(a\left(u+\frac{1}{u}\right)\right)\left(1-\frac{1}{u^2}\right) \dd u + O(N^{-1}).
\end{align*}

We only need to prove
\[
\frac{N}{a}\frac{1}{2\pi i}\oint\exp\left(a\left(u+\frac{1}{u}\right)\right)\dd u=N\int_{-2}^2\exp(ax)\frac{1}{2\pi}\sqrt{4-x^2}\dd x.
\]

We transform the contour to the unit circle and change variables $x=u+\frac{1}{u}$ ($u=\frac{1}{2}(x\pm i\sqrt{4-x^2}),\dd u=\frac{1}{2}(1\pm i\frac{-x}{\sqrt{4-x^2}})\dd x$)
\begin{align*}
\frac{a^{-1}}{2\pi i}\oint\exp\left(a\left(u+\frac{1}{u}\right)\right)\dd u&=\frac{a^{-1}}{2\pi i}\left(\int_{-2}
^2e^{ax}\frac{1}{2}(1+ i\frac{-x}{\sqrt{4-x^2}})\dd x+\int_{-2}
^2e^{ax}\frac{1}{2}(1- i\frac{-x}{\sqrt{4-x^2}})\dd x\right)\\
&=\frac{a^{-1}}{2\pi}\int_{-2}^2e^{ax}\frac{x}{\sqrt{4-x^2}}\dd x\overset{\text{by parts}}{=}\frac{a^{-1}}{2\pi}\int_{-2}^2(ae^{ax})\sqrt{4-x^2}\dd x
\end{align*}
\textbf{Conclusion:}
\[
\EE\left[\sum_{i=1}^N\exp\left(a\frac{\lambda_i}{\sqrt{N}}\right)\right]=N\int_{-2}^2e^{ax}\frac{1}{2\pi}\sqrt{4-x^2}\dd x+O\left(N^{-1}\right)
\]

\paragraph{Step 2: Variance}
\begin{align*}
    &\EE\left[ \sum_{i=1}^N \exp\left(a_1\frac{\lambda_i}{\sqrt{N}}\right)  \sum_{i=1}^N \exp\left(a_2\frac{\lambda_i}{\sqrt{N}}\right)\right] - \EE\left[ \sum_{i=1}^N \exp\left(a_1\frac{\lambda_i}{\sqrt{N}}\right)\right] \EE\left[ \sum_{i=1}^N \exp\left(a_2\frac{\lambda_i}{\sqrt{N}}\right)\right]\\
    =& N\frac{(a_1a_2)^{-1}}{(2\pi i)^2} \oiint \prod_{k=1}^2 \left( \frac{v_k + \frac{a_k}{\sqrt{N}}}{v_k}\right)^N \exp\left(\frac{a_k^2}{2N} + \frac{a_kv_k}{N}\right)\cdot \left[ \frac{v_1-v_2+\frac{a_1}{\sqrt{N}}-\frac{a_2}{\sqrt{N}}}{(v_1-v_2-\frac{a_2}{\sqrt{N}})(v_1-v_2+\frac{a_1}{\sqrt{N}})} - 1\right]\dd v_1 \dd v_2.
    \\
    \approx&N^2\frac{(a_1a_2)^{-1}}{2\pi i^2}\oiint\exp\left(a_1\left(v_1+\frac{1}{v_1}\right)+a_2\left(v_2+\frac{1}{v_2}\right)\right)
    \left[\frac{1+\frac{a_1-a_2}{N(u_1-u_2)}}{(1-\frac{a_2}{N(u_1-u_2)})(1+\frac{a_1}{N(u_1-u_2)})}\right]\dd u_1\dd u_2.
    \\
    =& N^2 \frac{(a_1a_2)^{-1}}{(2\pi i)^2} \oiint \exp\left(a_1\left(v_1+\frac{1}{v_1}\right)+a_2\left(v_2+\frac{1}{v_2}\right)\right)
    \left[\frac{a_1 a_2}{N^2(u_1-u_2)^2}\right]\dd u_1\dd u_2.
\end{align*}
\textbf{Conclusion:} Variance $\to \frac{1}{(2\pi i)^2} \oiint \exp\left(a_1\left(v_1+\frac{1}{v_1}\right)+a_2\left(v_2+\frac{1}{v_2}\right)\right) \frac{\dd u_1\dd u_2}{(u_1-u_2)^2}$.

Match this with the formula in Theorem~\ref{RMT_CLT}. Hint: Either directly with the 1st formula for $C(f,g)$, or
\[
\sum_{i=1}^Nf(\lambda_i)=\frac{1}{2\pi i}\oint_{\text{around all }\lambda_i}f(z)\sum_{i=1}^N\frac{1}{z-\lambda_i}\dd z
\]
and use it to compute with $C(\frac{1}{z-x},\frac{1}{w-x})$ in Theorems.

\paragraph{Step 3: Gaussianity} We need to show that 
\[
\EE\left[\prod_{k=1}^m
\left(\sum_{i=1}^N\exp\left(a_k\frac{\lambda_i}{\sqrt{N}}\right)-\EE\sum_{i=1}^N\exp\left(a_k\frac{\lambda_i}{\sqrt{N}}\right)\right)\right] \to \text{expressions of Wick's formula}.
\]
All of them have exactly the same $\prod_{k=1}^m$ part, but cross term part $\prod_{k<l}$ varies.
As in Steps 1,2, we change variables $u_k=\sqrt{N}v_k$ and use
\[
\text{crossterm}=1+\frac{a_ka_l}{N^2(u_k-u_l)^2}+O\left(N^{-3}\right)
\]
The summation over $2^m$ integrals leads to cancellation at parts involving 1.
The next term $\to$ perfect matching.
\end{proof}

\subsection{Spectrum Separation}
\todo

\subsection{Replica Method}
\todo
